\documentclass[12pt,a4paper]{article}

\usepackage[T2A]{fontenc}
\usepackage[utf8]{inputenc}
\usepackage[russian]{babel}
\usepackage{amsmath,amssymb,amsfonts}
\usepackage{booktabs}
\usepackage{longtable}
\usepackage{array}
\usepackage{geometry}
\usepackage{enumitem}
\usepackage{float}

\geometry{margin=2.2cm}

\title{Статические математические модели ИИС: версии 1--4}
\author{Исследовательский проект по интегральному индексу состояния}
\date{\today}

\begin{document}
\maketitle

\section{Назначение документа}
Настоящий документ фиксирует четыре статические версии модели ИИС (Интегрального Индекса Состояния), реализованные в проекте анализа физиологических данных. Под \textbf{статической} версией здесь понимается модель, которая рассчитывает значение индекса для одного сегмента данных независимо от соседних сегментов. Формальной отдельной динамической версии модели на момент составления документа не введено, поэтому в данном тексте рассматриваются только статические формулы.

\section{Общие обозначения}
Во всех версиях используются следующие обозначения:
\begin{itemize}[leftmargin=1.25cm]
    \item $A$ --- асимметрийный или латерализационный EEG-компонент.
    \item $\Gamma$ --- gamma-компонент корковой активации.
    \item $H$ --- гормональный или нейрогуморальный компонент.
    \item $V$ --- вегетативный компонент.
    \item $Q$ --- интегральный компонент качества/валентности состояния.
    \item $K$ --- коэффициент модуляции дофаминового вклада.
    \item $\sigma(x)=\dfrac{1}{1+e^{-x}}$ --- логистическая сигмоида.
    \item $\tanh(x)$ --- гиперболический тангенс.
    \item $\varepsilon$ --- малый стабилизирующий параметр.
\end{itemize}

Основные физиологические признаки:
\begin{itemize}[leftmargin=1.25cm]
    \item $\alpha_L,\alpha_R$ --- alpha-мощность слева и справа.
    \item $N_L,N_R$ --- общая мощность EEG по левому и правому полушарию.
    \item $P_\gamma$ --- мощность gamma-диапазона.
    \item $HF, LF$ --- компоненты вариабельности сердечного ритма.
    \item $DA$ --- дофамин или его proxy-оценка.
    \item $C$ --- кортизол или его proxy-оценка.
\end{itemize}

\section{Режимы расчёта}
Модели используются в трёх режимах:
\begin{enumerate}[leftmargin=1.25cm]
    \item \textbf{strict}: в расчёт допускаются только напрямую измеренные компоненты.
    \item \textbf{proxy}: допускаются proxy-оценки недостающих компонентов.
    \item \textbf{hybrid}: отсутствующие блоки исключаются, а веса оставшихся блоков автоматически перенормируются.
\end{enumerate}

\section{Версия 1}
\subsection{Смысл версии}
Версия 1 представляет собой раннюю четырёхкомпонентную структуру:
\[
IIS_1 = \sigma\!\left(0.35A + 0.25\Gamma + 0.30H + 0.10V\right).
\]

\subsection{Компоненты}
\paragraph{Компонент $A$.}
\[
A = \frac{\alpha_L - \alpha_R}{\alpha_L + \alpha_R + \varepsilon}.
\]
Этот компонент отражает alpha-асимметрию и связан с межполушарным распределением активности.

\paragraph{Компонент $\Gamma$.}
В текущей рабочей реализации используется калиброванная логарифмическая форма:
\[
\Gamma = \ln\!\left(1 + P_\gamma \cdot 10^{12}\right).
\]
Множитель $10^{12}$ переводит мощность из $V^2$ в $\mu V^2$, чтобы открытые EEG-данные не давали вырожденных почти нулевых значений.

\paragraph{Компонент $H$.}
\[
H = \frac{DA}{50} - \frac{C}{15}.
\]
На открытых датасетах прямые измерения $DA$ и $C$ отсутствуют, поэтому в proxy-режиме используются прокси по self-report и меткам состояний.

\paragraph{Компонент $V$.}
\[
V = \ln\!\left(1 + \frac{HF}{LF}\right).
\]
Логарифмизация введена для подавления редких выбросов HRV.

\paragraph{Компонент $Q$.}
В версии 1 компонент $Q$ не входит в итоговую формулу, но считается как диагностический:
\[
Q_{\mathrm{raw}} = \frac{\alpha_L - \alpha_R}{\alpha_L + \alpha_R + \varepsilon},
\qquad
Q = \tanh(\beta_q Q_{\mathrm{raw}}),
\]
где в реализации используется $\beta_q = 2.0$.

\subsection{Статус проверяемости}
Версия 1 полностью вычисляется на открытых EEG+ECG-датасетах в режиме proxy. На датасетах без EEG эта версия теряет основные различительные блоки.

\section{Версия 2}
\subsection{Смысл версии}
Версия 2 сохраняет ту же итоговую структуру, что и версия 1:
\[
IIS_2 = \sigma\!\left(0.35A + 0.25\Gamma + 0.30H + 0.10V\right),
\]
но меняет определение $A$.

\subsection{Компоненты}
\paragraph{Компонент $A$.}
\[
A = \frac{N_L - N_R}{N_L + N_R + \varepsilon}.
\]
Здесь вместо alpha-асимметрии используется асимметрия суммарной мощности по левому и правому полушарию.

\paragraph{Компоненты $\Gamma$, $H$, $V$, $Q$.}
Компоненты $\Gamma$, $H$, $V$ и диагностический $Q$ совпадают с версией 1.

\subsection{Статус проверяемости}
Версия 2 так же, как и версия 1, вычислима на открытых EEG+ECG-датасетах в режиме proxy. На реальных данных версия 2 показала почти те же разделительные свойства, что и версия 1.

\section{Версия 3}
\subsection{Смысл версии}
Версия 3 вводит дополнительный коэффициент модуляции $K$ и делает гормональный блок зависимым от него:
\[
IIS_3 = \sigma\!\left(0.35A + 0.25\Gamma + 0.30H + 0.10V\right),
\qquad
H = \frac{K \cdot DA}{50} - \frac{C}{15}.
\]

\subsection{Компоненты}
\paragraph{Компонент $A$.}
\[
A = \frac{N_L - N_R}{N_L + N_R + \varepsilon}.
\]

\paragraph{Компонент $\Gamma$.}
\[
\Gamma = \ln\!\left(1 + P_\gamma \cdot 10^{12}\right).
\]

\paragraph{Компонент $V$.}
\[
V = \ln\!\left(1 + \frac{HF}{LF}\right).
\]

\paragraph{Компонент $Q$.}
\[
Q_{\mathrm{raw}} = \frac{\alpha_L - \alpha_R}{\alpha_L + \alpha_R + \varepsilon},
\qquad
Q = \tanh(\beta_q Q_{\mathrm{raw}}),
\qquad
\beta_q = 2.0.
\]

\paragraph{Компонент $K$.}
В proxy-реализации на открытых данных:
\[
K = \frac{R \cdot G}{1 + \alpha S^2 + \beta F(S)},
\]
где
\[
F(S) = \frac{1}{1 + \exp(-k_f(S-S_0))}.
\]

В рабочей реализации:
\begin{itemize}[leftmargin=1.25cm]
    \item $R$ строится по HRV как ограниченная функция от $\dfrac{HF}{LF}$.
    \item $G$ строится по калиброванной gamma-мощности.
    \item $S$ строится как нормированная proxy-оценка стрессовой нагрузки.
    \item $\alpha = 1.2$, $\beta = 0.24$, $k_f=5.0$, $S_0=0.6$.
\end{itemize}

\paragraph{Компонент $H$.}
\[
H = \frac{K \cdot DA}{50} - \frac{C}{15}.
\]

\subsection{Статус проверяемости}
Версия 3 является проверяемой только в proxy-режиме, поскольку на открытых датасетах прямые измерения нейромедиаторов и гормонов отсутствуют. На реальных данных версия 3 показала устойчивое, но более слабое разделение состояний по сравнению с версиями 1 и 2.

\section{Версия 4}
\subsection{Смысл версии}
Версия 4 построена как \textbf{гормон-независимая} и ориентирована на эмоциональное состояние, вычислимое по реально доступным EEG/ECG/HRV-признакам. Формально:
\[
IIS_4 = \sigma\!\left(0.10A - 0.05\Gamma + 0.25V + 0.60Q\right).
\]

\subsection{Калибровочный принцип}
В версии 4 используется мягкая центрированная нормировка:
\[
T(x;c,w)=\tanh\!\left(\frac{x-c}{w}\right),
\]
где $c$ и $w$ подобраны по реальным диапазонам датасетов ds002722 и ds002724.

\subsection{Компоненты}
\paragraph{Компонент $A$.}
Версия 4 использует смесь alpha-асимметрии и общей межполушарной асимметрии:
\[
A_\alpha = T(\alpha_{\mathrm{asym}}; c_\alpha, w_\alpha),
\qquad
A_{\mathrm{total}} = T\!\left(\frac{N_L-N_R}{N_L+N_R+\varepsilon}; c_N, w_N\right),
\]
\[
A = \tanh\!\left(1.10A_\alpha + 0.40A_{\mathrm{total}}\right).
\]

\paragraph{Компонент $\Gamma$.}
Версия 4 трактует gamma как маркер избыточной активации:
\[
G_{\log} = T\!\left(\ln(1+P_\gamma \cdot 10^{12}); c_\gamma, w_\gamma\right),
\]
\[
G_{\gamma/\alpha} = T\!\left(\ln(1+\gamma/\alpha); c_{\gamma/\alpha}, w_{\gamma/\alpha}\right),
\]
\[
\Gamma = \tanh\!\left(0.30G_{\log} + 0.70G_{\gamma/\alpha}\right).
\]

\paragraph{Компонент $V$.}
Вегетативный блок в версии 4 смещён в сторону ЧСС как более устойчивой величины:
\[
V_{hr} = T(HR_0 - HR; 0, w_{hr}),
\qquad
V_{hrv} = T\!\left(\ln\!\left(1+\frac{HF}{LF}\right); c_{hf/lf}, w_{hf/lf}\right),
\]
\[
V = \tanh\!\left(0.80V_{hr} + 0.20V_{hrv}\right).
\]
Если HRV недоступна, используется только $V_{hr}$.

\paragraph{Компонент $Q$.}
Сначала рассчитывается физиологическое ядро:
\[
Q_{\mathrm{phys}} = \tanh\!\left(0.55A + 0.45V - 0.25\Gamma\right).
\]

В proxy-режиме оно смешивается с мягким якорем по self-report:
\[
Q_{\mathrm{anchor}}=
\tanh\!\left(
2.2\left(0.62V_n + 0.38(1-A_n) - 0.5\right)
\right),
\]
где $V_n$ --- нормированная valence, а $A_n$ --- нормированная arousal.

Итог:
\[
Q = (1-\lambda)Q_{\mathrm{phys}} + \lambda Q_{\mathrm{anchor}},
\qquad
\lambda = 0.35.
\]

\paragraph{Компоненты $H$ и $K$.}
В версии 4 отсутствуют:
\[
H \notin IIS_4,
\qquad
K \notin IIS_4.
\]

\subsection{Интерпретация}
Версия 4 является наиболее физиологически честной из четырёх версий в том смысле, что не использует трудноизмеримые гормональные переменные. Однако в proxy-режиме компонент $Q$ частично опирается на self-report, поэтому версия 4 пока должна трактоваться как \textbf{гибридная эмоциональная модель}, а не полностью независимый физиологический индекс.

\section{Эмпирическая проверка на реальных данных}
\subsection{Использованные датасеты}
Для практической проверки были использованы:
\begin{itemize}[leftmargin=1.25cm]
    \item \textbf{WESAD} --- полезен для сценария baseline/stress/amusement, но не содержит EEG.
    \item \textbf{CASE} --- полезен для динамики и continuous annotation, но не содержит EEG.
    \item \textbf{ds002722} --- содержит EEG, ECG и self-report.
    \item \textbf{ds002724} --- содержит EEG, ECG и self-report.
\end{itemize}

Для полноценного сравнения версий 1--4 наиболее важными оказались ds002722 и ds002724, поскольку только они дают одновременно EEG и ECG.

\subsection{Сравнительные результаты на ds002722}
\begin{longtable}{>{\raggedright\arraybackslash}p{2.6cm}cccc}
\toprule
Версия & Effect size & Overlap & Corr$(IIS, arousal)$ & Надёжность \\
\midrule
\endfirsthead
\toprule
Версия & Effect size & Overlap & Corr$(IIS, arousal)$ & Надёжность \\
\midrule
\endhead
IISVersion1 & $-2.410$ & $0.212$ & $-0.280$ & high \\
IISVersion2 & $-2.412$ & $0.224$ & $-0.281$ & high \\
IISVersion3 & $-1.853$ & $0.329$ & $-0.273$ & high \\
IISVersion4 & $-0.959$ & $0.522$ & $-0.199$ & medium \\
\bottomrule
\end{longtable}

\subsection{Сравнительные результаты на ds002724}
\begin{longtable}{>{\raggedright\arraybackslash}p{2.6cm}cccc}
\toprule
Версия & Effect size & Overlap & Corr$(IIS, arousal)$ & Надёжность \\
\midrule
\endfirsthead
\toprule
Версия & Effect size & Overlap & Corr$(IIS, arousal)$ & Надёжность \\
\midrule
\endhead
IISVersion1 & $-2.687$ & $0.174$ & $-0.180$ & high \\
IISVersion2 & $-2.693$ & $0.174$ & $-0.180$ & high \\
IISVersion3 & $-2.178$ & $0.273$ & $-0.185$ & high \\
IISVersion4 & $-1.446$ & $0.455$ & $-0.233$ & medium \\
\bottomrule
\end{longtable}

\section{Краткий итог по версиям}
\begin{enumerate}[leftmargin=1.25cm]
    \item \textbf{Версия 1} --- базовая, хорошо разделяет состояния на открытых EEG+ECG данных, но зависит от proxy-гормонов.
    \item \textbf{Версия 2} --- практически повторяет силу версии 1 и на ряде датасетов даёт лучший интегральный ранг.
    \item \textbf{Версия 3} --- интересна теоретически из-за коэффициента $K$, но на открытых данных уступает версиям 1 и 2 из-за сложности proxy-гормонального блока.
    \item \textbf{Версия 4} --- методически более чистая для эмоционального состояния, не использует гормоны, уже проверяема на реальных данных, но пока уступает по силе разделения proxy-усиленным версиям 1--3.
\end{enumerate}

\section{Вывод}
Четыре статические версии ИИС уже образуют полноценное семейство моделей, пригодных для сравнительного анализа на реальных физиологических данных. Версии 1--3 сильнее разделяют состояния в proxy-режиме за счёт дополнительного гормонального блока, тогда как версия 4 представляет собой более строгую и удобную для дальнейшего развития эмоциональную ветку модели, основанную только на реально доступных EEG/ECG/HRV-признаках.

\end{document}
